\begin{proposition}\label{prop:nonic57-exclusion}
An elliptic curve over a number field of degree $9$ has no point
of order $57$.
\end{proposition}
\begin{proof}
Put $X=X_1(57)$. Its genus is $85$, and it has good reduction at
$2$. A hypothetical level structure gives a noncuspidal closed
trace divisor $D$ of degree $9$, since order $57$ does not occur
in degree $1$ or $3$.
We use the partial reduction of
Proposition~\ref{prop:nonic57-quotient9-partial}.
It excludes noncuspidal degree-$7$ and degree-$9$ supports and
the two non-global nonic cusps at $2$.
Thus an entirely cuspidal reduction is globally intact.
Otherwise its noncuspidal part is a single degree-$6$ or
degree-$8$ place, with a globally rational cuspidal complement.
We treat these possibilities separately.

The exact winding calculation shows that the only simple factor
of $J_X$ not accounted for by the rank-zero winding quotient is
the elliptic factor $57\mathrm a1$.
Its Hecke eigenvalues at $2,5,7,11,13$ are
\[
 -2,-3,-5,1,2.
\]
In particular every diamond difference has finite rational image.
For the rational $q$-expansion, or dual, model put
\[
 \Phi_{n,\ell}=(T_n-a_n)
       (T_\ell-\langle\ell\rangle-\ell),\qquad
 n\in\{2,5,7,11,13\},\quad \ell\in\{5,7,11,13\}.
\]
Here $T_\ell$ denotes the geometric $\mu$-level quotient
correspondence, whose pullback is the retained canonical
$q$-Hecke matrix. Its reduction is
$\langle\ell\rangle\mathrm{Frob}_\ell+
\ell\mathrm{Frob}_\ell^{-1}$.
The first factor has finite image on $J_X(\mathbb Q)$.
The second kills rational torsion by Eichler--Shimura and
injectivity at the odd good prime $\ell$.
Hence every $\Phi_{n,\ell}$ annihilates $J_X(\mathbb Q)$.
Their pullbacks of differentials span a space of dimension $82$
over $\mathbb F_2$. The computation uses the full integral
canonical lattice and does not assume that $J_X(\mathbb Q)$
has rank zero.

We first exclude entirely cuspidal reduction.
The globally intact cusp places have signature $1^{18}2^9$.
For a rational cuspidal divisor $C$ lifting the proposed
reduction, apply the product of the maps $\Phi_{n,\ell}$ to
$[D-C]$. Its value is zero.
It suffices to prove formal immersion on $X^{(9)}$ at the
reduction of $C$.

We obtain the required cusp parameters directly.
A saturated canonical basis is retained through $q^{1100}$.
The rational cusp parameter is its usual expansion divided by
$q$. For a quadratic cusp use the genuine Tate family
\[
 E(q^{19}),\qquad P=\zeta_3q.
\]
The two modular units $x,y$ have pole degrees $37,39$,
verified independently at all $72$ geometric cusps.
The $49$ and $47$ canonical unit identities are global:
the corresponding section degrees are at most $205,207$,
whereas their verified orders are at least $1098$.
They recover the quadratic cusp arc through order $20$.
An independent valuation-pivot calculation gives valuation $95$
for the chosen integral $84\times84$ minor and verifies every
coefficient of this arc. Its canonical tangent is nonzero.
The diamond operator generates the $18$ rational cusps and the
$9$ closed quadratic cusps. Its ninth power on the latter arc
is coefficientwise Frobenius, as checked explicitly.

Every degree-$9$ cuspidal divisor has a rational support point.
Translate one of maximal rational multiplicity to the first cusp.
If that multiplicity is $n_0$, the other rational multiplicities
are at most $\min(n_0,4)$, and a quadratic multiplicity is at most
$4$. We test the corresponding differential prefixes, using both
$\mathbb F_2$-trace components at every quadratic cusp.
There are exactly
\[
 265729,370481,122514,29818,6267,1122,162,17,1
\]
normalized configurations for $n_0=1,\ldots,9$.
All $796111$ matrices have rank $9$.
This is an independent complete enumeration; its counts also
agree with the respective generating-function coefficients.
Thus the product map is a formal immersion at every proposed
$C$. Since $D$ and $C$ have the same reduction and the same image,
formal immersion gives $D=C$, a contradiction.

It remains to treat noncuspidal reduction.
We now use the point-level geometric quotient correspondence
$(E,P)\mapsto(E/G,\phi(P))$. Its reduction is
$\mathrm{Frob}_5+5\sigma\mathrm{Frob}_5^{-1}$.
We put
$\sigma=\langle5\rangle$ and $A_5=T_5-1-5\sigma$.
The class $[\sigma D-D]$ is rational torsion: its normalized
class has finite diamond-difference image, and the remaining
cusp difference is torsion by Manin--Drinfeld.
Thus $A_5[\sigma D-D]=0$.
For the cuspidal complement $C$, the class $[\sigma C-C]$
is rational cuspidal torsion, since $C$ is globally rational.
It is therefore also annihilated by $A_5$.
Consequently, if $P$ is the noncuspidal residue place and
$T=T_5P$, the degree-zero divisor
\[
 B=\sigma T-T+P+4\sigma P-5\sigma^2P
\]
must be principal.

Complete twist enumeration leaves two supersingular elliptic
models over $\mathbb F_{64}$, both of order $57$, and one
Galois orbit of ordinary models over $\mathbb F_{256}$, of order
$228$. Their $57$-primary groups are cyclic, so their level
structures are covered by diamond translates of one point on
each retained model. We check both supersingular representatives
and one ordinary representative. Hecke operators and $\sigma$
commute with every diamond, so these checks cover all levels.

Every cyclic $5$-kernel and its actual level image is retained.
For either degree-$6$ representative, $T_5P$ has closed degrees
$18,9$ with norm multiplicities $1,2$.
For the degree-$8$ representative the closed degrees are $24,24$,
each with multiplicity $1$.
Independent isogeny calculations verify the kernel polynomials,
all actual point images and all norm multiplicities; the products
of their norms are exactly the respective $5$-division polynomials.
The points $P,5P,25P$ and both target levels are computed directly,
so no unverified diamond convention enters $B$.
Equal closed places are identified before adding signed
multiplicities. The positive and negative degrees of $B$ are
$42,42,88$ for the three retained representatives.

We recover sections of $2K_X$ using its $215$ and $213$ unit
identities. Their section degrees are at most $373,375$ and are
again below the verified expansion orders.
For an ordinary model we deform $a_6$ to $a_6+t$.
For a supersingular model, where $a_1=0$ and $a_3\ne0$, we
instead deform $a_1$ to $t$.
The general $57$-division recurrence lifts the level point
uniquely; its $x$-derivative and the elliptic equation's
$y$-derivative are nonzero.
The elliptic equation, division recurrence and unit identities
are checked through order $16$.
Independent recovery of the canonical coordinates verifies the
same $2K_X$ arc and a nonzero tangent. These are therefore genuine
local parameters, including at the supersingular points.

Put $R_i=H^0(X,iK_X)$. Products of the canonical basis give
$\dim R_2=252$, and products of this basis give $\dim R_4=588$.
All $3655$ and $31878$ unordered product identities are checked
through $q^{1100}$, beyond the respective Sturm bounds $480,960$.
For $B=B_+-B_-$ choose nonzero $F,H\in R_2$ vanishing on $B_+$.
If $f\in\mathcal L(B)$, then $G=fF$ vanishes on $B_-$ and
$GH=F(fH)\in FR_2$.
It therefore suffices to verify injectivity of
\[
 R_2\longrightarrow
 H^0(B_-,2K_X|_{B_-})\oplus R_4/FR_2,
 \qquad G\longmapsto(G|_{B_-},GH\bmod FR_2).
\]
Multiplication by $F$ has rank $252$.
The vanishing-condition ranks are exactly $42,42,88$.
All three final matrices have rank $252$, and one section $H$
suffices in each case. Independent Sage calculations reconstruct
the entire canonical multiplication table and every final matrix.
Thus $\mathcal L(B)=0$ for every remaining support, contradicting
principality. This excludes order $57$.
\end{proof}
